\documentclass[sigconf]{acmart}

% ============================================================
% Draft settings: hide ACM copyright/reference block while the
% outline is under review. Remove these four lines for the real
% SIGCOMM submission and follow that year's Call for Papers.
% ============================================================
\settopmatter{printacmref=false}
\setcopyright{none}
\renewcommand\footnotetextcopyrightpermission[1]{}
\pagestyle{plain}
\raggedbottom

\usepackage{booktabs}
\usepackage{graphicx}
\usepackage{xcolor}


\begin{document}

\title{LIMER: Link-Level Detection and Localized Recovery
in Active-Active GPU Training Fabrics}

\author{Author One}
\authornote{Author list and affiliations to be finalized.}
\affiliation{%
  \institution{Institution}
  \country{}}
\email{first.author@example.com}

\begin{abstract}
Gray link failures preserve connectivity while impairing communication.
We study how to confirm an affected path below 100~ms in an active-active
dual-plane fabric, with explicit observation cost, and how that decision
could support beneficial, coordinated local rerouting. The primary
detection candidate separates business and probe observations, tracks
guarded receiver sequence gaps and compensates for sparse business traffic
with representative probes. In a registered 10-Gb/s one-hop packet model,
fresh validation gives 400/400 correctly localized decisions below 100~ms,
a largest per-condition p95 of 75.4~ms and no alarms in 60 primary controls.
Another 60 high-traffic controls are alarm-free. Average probe traffic is
10.84--299.04~Mb/s/path across the tested business profiles. Same-report
comparisons distinguish decision-rule effects from observation opportunity,
and smaller-budget and nonrepresentative-probe failures remain visible.

Separate component experiments study the use of an impairment alarm.
A bounded ACK-progress route loop reduces completion time by 13.34\%
with a faster alternative and avoids the 35.42\% penalty of fixed switching
with a slower alternative. Versioned pending-flow checks preserve ownership
under modeled coordination. An independently reconstructed OptCC comparator
provides the scheduling baseline for future joint comparisons. The primary
detector has not been integrated into the collective route loop. These
results establish conditional simulation behavior, not a universal timing
guarantee or a validated end-to-end recovery system.
\end{abstract}

\keywords{GPU clusters, collective communication, fault tolerance,
gray failures, datacenter networks, active-active fabrics}
\maketitle

\section{Introduction}
An access link can remain UP while packet loss or severe delay slows the
communication that depends on it. In a collective, a delayed transfer can
hold back later dependent work. Our primary question is whether a
communication-impaired path can be correctly confirmed strictly below
100~ms of fault onset. We require at least 95\% timely decisions in each
registered eligible fault condition, rather than a favorable average across
conditions. Late, missed and unverified attempts remain in the denominator.

Fast confirmation requires observations that expose the impairment seen by
business traffic. Passive transport signals can arrive late or respond to
reordering and congestion. Active probes create observations when business
traffic is sparse, but a small probe can remain healthy while larger business
packets suffer. More probe traffic can also improve detection without making
the decision rule better. Detection speed, observation representativeness
and observation cost must therefore be evaluated together.

LIMER's primary detection candidate separates business and probe acceptance,
delivery and sequence observations. Receiver gaps must persist through a
guard before they supply impairment evidence. Representative probes follow
the same modeled delivery path as business packets; their source rate uses
only the previous closed business-traffic window. This makes the extra
observation budget causal and measurable. It does not guarantee that probes
experience every fault that business traffic might encounter.

An impairment decision is useful only if the subsequent action helps.
Both network planes already carry work in our active-active setting, so
rerouting consumes service that is also needed by the alternate plane's
existing work. A route decision must compare the remaining work and
coordination cost. Once chosen, a local update also needs a shared version
and activation point so that senders and receivers agree. These are
supporting questions around the primary detection constraint, not an
assumption that every alarm should cause a switch.

We organize the study around three questions:
\begin{enumerate}
\item \textbf{Fast detection:} which observations allow at least 95\% of
decisions to be timely and correct in each fault condition, and at what cost?
\item \textbf{Route utility:} when does localized switching reduce the
completion time of finite remaining communication compared with staying?
\item \textbf{Execution consistency:} which coordination and transfer-ownership
checks are required to activate a route change within unfinished work?
\end{enumerate}

The strongest detection validation has 400/400 timely fault decisions,
largest per-condition p95 75.4~ms and 60 alarm-free primary controls in an
explicit 10-Gb/s one-hop model. A separate high-traffic supplement has
60 alarm-free controls. Under adaptive observations, the same-report shadow
gives 391/400 timely decisions and misses the idle condition's target;
under a constant high probe budget, both rules meet each condition's target.
This distinction limits the attribution of gains to the decision rule.

The recovery evidence is separate. An online loop driven by ACK progress
demonstrates a beneficial switch with faster alternate service and a stay
decision with slower service. Checks of pending transfers validate
coordination and ownership under a modeled reliable management channel.
The primary packet detector and its observation traffic have not been
integrated into that collective loop. Our results therefore support a
conditional detection finding and supporting recovery components, not a
completed recovery system. The study ends at simulation and
CPU emulation; physical deployment is outside its scope.

\section{Problem Setting and Comparison Scope}
\label{sec:setting}
\subsection{Active-Active Communication and Impairment}
A plane is an independently routed network connecting the participating
workers. Both planes carry communication slots in the initial plan, so the
alternate plane is already active. A slot identifies a sender and a
communication step. Local rerouting changes affected slots while retaining
the workers and communication schedule. A ring schedule has $2(p-1)$
reduce-scatter/all-gather communication steps for $p$ ranks, with endpoint
dependencies that let a slow transfer delay subsequent work.

A hard fault removes carrier; a gray fault preserves carrier while
impairing delivery. We distinguish random packet loss, configured service-rate
reduction and severe delivery delay. Their effects on transport are different.
The target decision concerns impaired delivery, including severe slowdown;
it need not establish permanent loss or its physical cause. The one-hop
detector labels identify sender/plane paths, not arbitrary physical hops
inside a multi-hop fabric.

\subsection{Timing Target and Evidence Boundaries}
Detection latency runs from fault onset to the delivered-report-driven
decision. In packet simulation, onset is the simulated injection timestamp;
in emulation, it is completion of the injection action. Each campaign
registers eligible severity and duration and retains every attempt. The
primary empirical gate requires at least 95\% timely, correctly localized
decisions in each eligible condition and no counted control or off-target
alarms. This finite-sample gate is not a population confidence guarantee.
Earlier all-trial and loss-only gates keep their original verdicts.

Detection, route commit and completed communication are different milestones.
We report their respective timings without adding measurements from separate
experiments into an end-to-end claim. Hard-fault detection below 1~ms and
progress resumed within 1~s of detection remain additional design targets,
not demonstrated global guarantees. Numerical AllReduce correctness is also
unvalidated: the workloads execute communication dependencies and framed
payloads, without performing tensor reduction.

\subsection{Related Work and Comparator Roles}
OptCC schedules AllReduce under post-failover bandwidth
asymmetry~\cite{optcc2026}. It selects a collective schedule, whereas our
route policy selects placement of affected work across active planes. These
decisions can be composed. We reconstruct OptCC as an executable scheduling
comparator and compare it with an ordinary ring; we do not equate OptCC
with staying on a damaged path. The crossed ring/OptCC and route-policy
comparison remains unrun.

Chameleon selects data rerouting or dynamic parallelism after node
failures~\cite{chameleon2026}, so adaptive recovery selection alone is
established prior work. SprayCheck uses per-spine distributions under packet
spraying to detect gray failures~\cite{spraycheck2026}. Its evaluation focuses
on spine faults and sketches an access-link count/NACK extension. It is a
mechanism reference for our slot-based access observer, not a reproduced
head-to-head detector. We do not convert iteration-based detection into
a claimed 100-ms comparison.

VCCL uses primary and backup queue pairs for port-failure recovery within
its collective library~\cite{vccl2026}. Varuna distinguishes unexecuted
operations from operations whose acknowledgments were lost, using completion
logs, extended atomic status and dynamic connected queue
pairs~\cite{varuna2026}. Its configured backup order is different from
our remaining-work cost heuristic. Its RDMA microbenchmark and Motor
TPC-C results are not training measurements in this study, and its logging
corner case and separate CAS argument do not prove TCP replay correctness.
These systems motivate explicit execution identity. They are literature
comparisons, not reproduced recovery baselines here.

\section{Fast Detection with Representative Observations}
\label{sec:detection}
\subsection{Observation Contract}
The primary candidate detects delivery impairment from observations
available at the two ends of each monitored path. Business packets and
probes have separate accepted/delivered counters and independent source
sequence numbers. Reports carry path, batch, traffic class, sample time,
arrival time and receiver gap evidence. Injection time, affected-path truth,
private drop counters and future traffic never enter the decision interface.
The collector runs the rule on report arrival, so reporting delay contributes
to the measured confirmation time.

Business and probe classes remain separate because successful probes can
otherwise hide business impairment. The receiver starts a sequence gap only
after a higher sequence arrives and clears it when the missing packet is
later delivered. A gap therefore records incomplete delivery rather than
proving permanent packet loss. The observation identifies a path in the
declared one-hop model.

\subsection{Guarded Impairment Decision}
A receiver gap persisting for at least 40~ms supplies impairment evidence.
Its age is computed at sampling, without counting report transit as extra
receiver evidence. Two aligned business-only service windows also detect
a sustained delivery decline that healthy probes could mask. The guarded
service rule requires sufficient accepted traffic and a delivered increment
no greater than 60\% of the accepted increment in two consecutive windows.
The aligned early lane uses at least 32 accepted packets; missing, invalid
or nonsevere windows break evidence.

The candidate retains the frozen hybrid observer's mature batch-deficit
and aligned total-service lanes as fallbacks. Their exact guard, window
and reset rules are described in Appendix~\ref{app:observers}. Only the
first alarm per path is emitted in a lifecycle. The frozen hybrid rule
is also scored independently on the same delivered reports as a shadow
baseline. This shared input permits rule attribution within a campaign;
results from different traffic models are not treated as a matched ranking.

\subsection{Representative Probes and Causal Compensation}
Representative probes use the business packets' 1500-byte wire size and
the same modeled socket, queue, route and application delivery stage.
This common fault treatment is an explicit assumption. Same-path small
probes do not imply equivalent impairment exposure: the matched selective-loss
contrast in Section~\ref{sec:attribution} demonstrates the distinction.

Sparse or idle business traffic provides too few observations, so a source
policy adds probes according to accepted traffic in the previous closed
10-ms business window. Let $D_{b-1}$ be that count and $\Delta_b$ the next
source span. The source evenly schedules
\[
P_b=\max\!\left(0,
\left\lfloor\frac{25000\Delta_b}{1\mathrm{s}}\right\rfloor-
\left\lfloor\frac{D_{b-1}\Delta_b}{10\mathrm{ms}}\right\rfloor\right)
\]
probes in that span. Startup from 1 to 10~ms assumes no prior business.
This policy uses neither receiver loss nor fault metadata. Historical
traffic can cause transient overshoot or lag, and every actual emitted
probe byte is charged. A constant 35,000-pps representative source is
evaluated separately to expose the effect of observation budget.

\subsection{Reporting Cost and Output}
Source and receiver clocks are aligned in the evaluated model. Reports
occupy 96 bytes per copy: 68 payload bytes plus modeled IPv4/UDP overhead,
without link-layer overhead on the abstract collector. A separate serialized
10-Mb/s collector adds 1-ms propagation. All 4800 offered report copies per
lifecycle cost 460.8~KB, with 7.3728~Mb/s steady offered load. This collector
does not share congestion with the monitored data paths. Its capacity and
delay are part of the observation envelope. Callback computation time is
not modeled as elapsed simulation time.

The detector outputs an impaired path and decision time for subsequent
route-cost evaluation. No route action follows from the detection experiment
itself. Multiple-episode rearming and arbitrary report loss, management delay
or unknown probe treatment remain outside the latest validation.

\section{Route Decisions and Coordinated Execution}
\label{sec:recovery-design}
The following components study how an impairment decision can be used.
They are evaluated separately from the primary detector. The bounded online
route loop uses native sender ACK progress, not the representative packet
detector in Section~\ref{sec:detection}. This distinction is retained in
every performance interpretation.

\subsection{Remaining-Work Switch-versus-Stay Heuristic}
Let $I_A,I_B$ be launched, unacknowledged bytes and $W_A,W_B$ pending
bytes on the two planes. Only pending affected A work can change routes.
The current observed ACK service is $r_A$; alternate service $r_B$ is the
minimum of its calibrated service and its last positive fully busy sample.
Unknown, zero or more-than-20-ms-old alternate service selects stay.
Otherwise the sender-service heuristic compares
\begin{align*}
\widehat T_{\rm stay}&=\max\{(I_A+W_A)/r_A,(I_B+W_B)/r_B\},\\
\widehat T_{\rm switch}&=\max\{I_A/r_A,(I_B+W_B+W_A)/r_B\}+c.
\end{align*}
The registered cost budget is $c=10~\mu$s, exceeding the default modeled
transaction cost. A switch requires pending A work and a smaller switch
estimate. Remaining work on B is charged because both planes are active.
The estimate is a sender-service proxy, not a collective critical-path
optimizer or a prediction of future alternate congestion.

For this component experiment, native sender ACK credit combines completed
flow bytes and active acknowledged bytes. Samples are taken every 2~ms and
reports pass through the modeled collection channel. A detector calibrates four positive fully
busy samples and alarms after three busy samples below half their reference;
idle windows interrupt evidence. These configured-rate development checks
do not establish the primary detector's broader weak-loss/idle coverage.
Appendix~\ref{app:online} retains the full detector and earlier policy details.

\subsection{Versioned Activation and Transfer Ownership}
At transfer launch, the SimAI target selects the committed plane, posts
one matching receive and retains that transfer's ports and version through
completion. A route update changes only pending affected transfers.
PREPARE, READY and COMMIT messages coordinate a common version and
activation point across all four ranks. The reliable modeled management
FIFO charges 64 bytes per message with 1-Gb/s serialization and 1-$\mu$s
propagation by default. Its messages are not packetized on the RDMA data
network. A missing READY prevents activation; a nonconsecutive version is
rejected.

Launched transfers retain ownership on their original paths. Consequently,
the pending-only mechanism cannot rescue a severely stalled launched
transfer. The separate CPU/TCP harness checks frame round, step, route
version, payload length and checksum, with explicit ownership of partially
read frames. These transport checks are documented in
Appendix~\ref{app:execution}; they are not an exactly-once numerical reduction
proof and do not imply transparent migration of in-flight work.

\section{Evaluation}
\subsection{Scope, Metrics and Baselines}
We evaluate fast detection first, then its observation cost and attribution,
and finally the separate route-utility and execution components. Packet-level
experiments use the pinned SimAI network frontend~\cite{simai2025} and ns-3 backend, with
separate observation traffic and management models. CPU/Mininet framed-TCP
emulation supplies additional implementation and failure evidence; its
wall-clock times are not pooled with simulated timings.
Appendix~\ref{app:implementation} records the implementation environments.

The primary detection metric is the fraction of all registered attempts
correctly localized strictly below 100~ms in each condition. We also report
late and missed decisions, per-condition p95, control/off-target alarms and
actual probe/report bytes. Source, parameters and dependencies are frozen
before fresh validation, and independent event replay checks the recorded
decisions. Development outcomes are not pooled into acceptance.

Detection comparators include a constant representative-probe source and
the frozen hybrid rule on the same delivered reports. Earlier TCP and
batch-only rules establish specific failure boundaries, not a matched
cross-platform leaderboard. Route arms are stay, fixed localized switching
requested at startup, oracle-timed switching supplied with fault time and
identity, and the observable candidate. Oracle timing does not provide
optimal route choice. OptCC and ring compare scheduling in a separate
network-only setting; no completed experiment compares full LIMER with OptCC.

\subsection{Fast Detection in the Registered Envelope}
\label{sec:primary-validation}
The primary validation contains eight independent one-hop UDP paths with
10-Gb/s links, a fixed affected path, 250-ms source activity and observation
drain through 600~ms. Six weak-fault conditions use 0.5\% persistent loss:
selective network/application loss at 80~Mb/s and 1~Gb/s business, sparse
1-Mb/s selective loss, and idle size-independent loss. Two severe-delay
conditions add 80~ms to 25\% of network or application deliveries. All paths
are monitored; fault time and identity enter scoring only.

After 130/130 timely development faults and 30 alarm-free development
controls, unchanged-source fresh validation produces 400/400 timely fault
decisions (Table~\ref{tab:adaptive-business-validation}). Every condition
meets the empirical target, with no late or missed decisions and largest
per-condition p95 75.4~ms. The 60 primary controls and a separately registered
60-case high-traffic supplement have no alarms. Controls include source
pauses and 40-ms network/application deferrals; these finite checks do not
establish a population false-alarm rate.

\begin{table}[t]
\caption{Causal traffic-compensation validation on 10-Gb/s links.
The detector and 40-ms guard are unchanged; each attempt stays in its denominator.}
\label{tab:adaptive-business-validation}
\centering
\small
\setlength{\tabcolsep}{3pt}
\begin{tabular}{@{}lrrrr@{}}
\toprule
Condition & $<100$ ms & Late & Miss & P95 ms \\
\midrule
Selective, 80 Mb/s & 60/60 & 0 & 0 & 70.1 \\
Application loss, 80 Mb/s & 60/60 & 0 & 0 & 72.1 \\
Selective, 1 Mb/s & 60/60 & 0 & 0 & 73.6 \\
Business idle & 60/60 & 0 & 0 & 75.4 \\
Selective, 1 Gb/s & 60/60 & 0 & 0 & 58.7 \\
Application loss, 1 Gb/s & 60/60 & 0 & 0 & 57.5 \\
Network delay 80 ms & 20/20 & 0 & 0 & 54.5 \\
Application delay 80 ms & 20/20 & 0 & 0 & 54.3 \\
\bottomrule
\end{tabular}
\end{table}

Independent replay reconstructs counters, gap ages, probe emissions,
encoded reports, collector service and native/Python alarm agreement.
The result is conditional on aligned clocks, representative fault treatment
and the stated reporting capacity. Even 60/60 and 20/20 successes have
two-sided Wilson 95\% lower bounds of only 93.98\% and 83.89\%, respectively.
An empirical pass therefore does not establish a population 95\% guarantee
or a collective recovery benefit.

\subsection{Observation Cost, Attribution and Retained Failures}
\label{sec:attribution}
Causal traffic compensation reduces extra observations when business traffic
already supplies them. Table~\ref{tab:probe-cost} charges actual accepted
probe bytes over the full 249-ms source interval, including startup and lag.
At 1-Gb/s continuous business, only 225 startup probes/path are emitted;
sparse and idle business still require substantial active sensing. The
separate 7.3728-Mb/s reporting load is additional to these probe costs.

\begin{table}[t]
\caption{Actual average adaptive probe traffic on each 10-Gb/s path.
The denominator includes the full source lifecycle, not only steady state.}
\label{tab:probe-cost}
\centering\small
\begin{tabular}{@{}lrr@{}}
\toprule
Business profile & Probes (Mb/s) & Link capacity (\%)\\
\midrule
Continuous 1 Gb/s & 10.84 & 0.108\\
Continuous 80 Mb/s & 223.18 & 2.232\\
Sparse 1 Mb/s & 299.04 & 2.990\\
Business idle scene & 274.60 & 2.746\\
\bottomrule
\end{tabular}
\end{table}

Same-report comparisons distinguish observation policy from decision-rule
effects. Under constant 1500-byte probes at 35,000 pps, nominal probe cost is
420~Mb/s/path, or 4.2\% of link capacity. Candidate validation gives 400/400
timely faults, while the frozen hybrid shadow gives 396/400 and also meets
every condition's 95\% gate. Most of that target pass therefore reflects the
changed observation opportunity. Under adaptive observations, candidate
validation again gives 400/400 while its same-report shadow gives 391/400;
idle coverage is 60/60 versus 55/60. The guarded gap rule supplies a
target-relevant difference in that adaptive campaign. Constant and adaptive
campaigns use fresh ledgers; they are not paired-seed algorithm-only arms.

Probe representativeness and budget remain necessary conditions.
An earlier matched 80-Mb/s contrast has the same business send/loss counts
and 123 dropped business packets in total in each arm: the small-probe
hybrid detects 20/20 IID-loss cases but only 5/20 size-selective cases while
its probes stay healthy. At the original 100-Mb/s, 17.92-Mb/s probe envelope,
class-separated development retains 16/20 timely selective and 13/20
application-loss decisions. Making probes business-shaped at the old byte
budget reaches 20/20 selective decisions but only 18/20 application-loss
and 2/5 idle decisions. These low-budget conditions retain their original
FAIL verdicts. Appendix~\ref{app:detection-evidence} retains
all earlier detection tables, guard sensitivity, controls and setup failures.

\subsection{When Localized Switching Helps}
\label{sec:route-results}
These experiments evaluate route utility separately from the primary
detector. A static four-rank, 3-MiB ring calibration fixes both access
rates across stay/switch arms and changes only the affected rank's A
sending slots. At 40/200~Gb/s, switching reduces completion from 0.740
to 0.571~ms (22.9\%); at 40/20~Gb/s, it increases completion from 1.454
to 1.934~ms (33.0\%). Table~\ref{tab:dual-calibration} excludes runtime
detection, coordination and replay costs. It motivates the route decision
but is not a controller result.

\begin{table}[t]
  \caption{Static dual-plane network calibration. Both route arms share
  each row's capacities; control, cutover and replay costs are excluded.}
  \label{tab:dual-calibration}
  \centering
  \small
  \begin{tabular}{@{}rrrr@{}}
    \toprule
    $C_A$ (Gb/s) & $C_B$ (Gb/s) & Stay (ms) & Switch (ms) \\
    \midrule
    100 & 200 & 0.352 & 0.305 \\
    40 & 200 & 0.740 & 0.571 \\
    40 & 20 & 1.454 & 1.934 \\
    \bottomrule
  \end{tabular}
\end{table}

The bounded online comparison uses the ACK-progress detector and
remaining-work heuristic of Section~\ref{sec:recovery-design}. Four ranks
execute fourteen 3-MiB rings with prescribed releases at
0,100,\ldots,1300~ms. At 312~ms, rank-2 A configured egress service falls
from 1 to 0.2~Gb/s; B initially supplies either 1 or 0.1~Gb/s. Other
accesses remain at 1~Gb/s. All four arms share workload and per-report
cost, with 96-byte path reports offering 3.072~Mb/s to the 10-Mb/s
collector. Each activated plan charges 768 modeled control bytes.

\begin{table}[t]
\caption{Bounded online ring comparison: total network-DAG completion
from workload start (ms). Fixed requests localized switching at startup;
oracle supplies fault timing/identity, not optimal route choice. All arms
share workload, reporting model and charged coordination.}
\label{tab:online-route}
\centering\small
\begin{tabular}{@{}rrrrr@{}}
\toprule
$C_B$ (Gb/s) & Stay & Fixed & Oracle & Candidate\\
\midrule
1.0 & 1545.674 & 1339.485 & 1339.485 & 1339.485\\
0.1 & 3958.544 & 5360.714 & 5190.811 & 3958.544\\
\bottomrule
\end{tabular}
\end{table}

All eight registered workloads complete. With faster B, the candidate
alarms after 9.2304~ms, changes 31 pending transfers and completes in
1339.485~ms rather than stay's 1545.674~ms, a 13.34\% reduction. With
slower B, it alarms after 35.2304~ms and stays, completing in 3958.544~ms;
fixed switching instead takes 5360.714~ms, a 35.42\% increase over stay.
Neither candidate case has a pre-event or unaffected-path alarm.
Independent replay audits 96,224 delivered reports, 2.114~GB of ACK
credit, policy arithmetic, dependencies and immutable launched ownership.

These deterministic configured-rate examples establish bounded route
selection behavior. They do not establish weak-loss coverage, optimality,
future alternate-capacity guarantees or performance with the primary
representative detector. Appendix~\ref{app:cost-diagnostics} retains the
static setup corrections and severe-loss TCP contrast.

\subsection{Execution Consistency Checks}
\label{sec:execution-results}
The scripted SimAI transaction gate has nine prospectively registered
engineering cases: eight workloads complete with transport-map drain and
one deliberately nonconsecutive version is rejected. The default request
at 5~$\mu$s activates at 13.145~$\mu$s, moving two still-pending affected
transfers while launched work retains its endpoints and version. Missing
READY prevents activation. With slower coordination, the workload can
finish before activation, so a valid update moves no transfers.

Independent reconstruction checks all 582 trace records and the frozen
dependencies. Separate framed-TCP runs check integrity and version
agreement, but retain transient false commits and severe-loss timeouts.
Appendix~\ref{app:execution} gives those failures and the complete transaction
cases. These checks validate the declared execution contract, not numerical
reduction, cancellation/replay or rescue of already-launched stalls.

\subsection{Executable Scheduling Comparator}
\label{sec:baseline-reproduction}
OptCC and ordinary ring schedules share the same network runner using
SimAI and ns-3. The matched setting has one nonblocking switch and one
participating GPU per server. Healthy NICs supply 200~Gb/s; the last-rank
NIC supplies 100~Gb/s ($L=2$). Both arms use 500-ns propagation,
9000-byte payload MTU, DCQCN and 3-$\mu$s launch delay. The setting
includes no packet loss, NVLink or PXN. OptCC uses
$k=4$, and both schedules transmit $2(p-1)N$ data bytes in aggregate;
OptCC synchronization nodes additionally carry 20-byte control payloads.

\begin{table}[t]
  \caption{Independent network-only replay, $L=2$, $k=4$. Ratio is ring
  makespan divided by OptCC makespan. These are reconstructed measurements,
  not reproduced paper figures or LIMER results.}
  \label{tab:optcc-replay}
  \centering
  \small
  \begin{tabular}{@{}rrrrr@{}}
    \toprule
    $p$ & $N$ (MiB) & Ring (ms) & OptCC (ms) & Ratio \\
    \midrule
    5 & 5 & 0.712 & 0.638 & 1.116 \\
    5 & 20 & 2.754 & 2.283 & 1.206 \\
    16 & 15 & 2.500 & 1.986 & 1.259 \\
    16 & 60 & 9.718 & 6.917 & 1.405 \\
    \bottomrule
  \end{tabular}
\end{table}

All 24 primary launches pass; three deterministic launches per arm and
condition reproduce the same values. OptCC reduces network makespan by
10.4--28.8\% relative to ring in these four settings. Four healthy-ring
references and two segmentation checks bring the measurement denominator
to 30/30. These repeated fixed settings do not provide confidence intervals.

The public artifact lacks the original simulator integration and experiment
driver, so this is an independent reconstruction, not an exact reproduction
of a published performance figure. It omits native collective memory and
reduction costs. Appendix~\ref{app:optcc} preserves artifact, transport,
scheduler and segmentation checks. The result establishes an executable
schedule comparator, not LIMER superiority. The planned common-backend
matrix crosses ring/OptCC schedules with stay/fixed-switch/LIMER route
policies under identical faults, workload and seeds, charging all recovery
cost and retaining failures. That joint comparison remains unrun.

\section{Limitations and Remaining Integration}
\label{sec:limitations}
The primary unresolved task is joint simulation: the representative
detector, its actual observation traffic and the dynamic collective
controller have not been evaluated together. Passing isolated detection
and component route tests does not establish timely or beneficial recovery
of a collective. The shorter online detector timings and the 400/400 packet
validation cannot be combined into one measured end-to-end result.

Detection coverage remains conditional. The latest model has aligned clocks,
eight independent one-hop paths, one fixed affected path and configured
collector service. Probe size/path/application treatment must expose the
business impairment. Smaller-budget selective/application/idle failures
remain visible, as do earlier reordering false alarms and delayed-report
deadline failures. The latest primary validation does not extend its pass
to arbitrary report loss, slow management, correlated faults, simultaneous
traffic changes or multiple fault episodes. Synchronous callback computation
also omits detector processing latency. General fast gray-fault coverage
therefore remains unresolved.

Route selection assumes observations can estimate useful alternate service
over finite remaining work. The current heuristic cannot optimize the
collective critical path or guarantee future service under changing
background traffic. Pending-only execution leaves launched stalls on their
old routes; unreliable coordination and numerical reduction require further
simulation validation. The current bounded ring result does not replace the
crossed ring/OptCC route comparison. Expansion to 128 simulated ranks remains
unrun; Appendix~\ref{app:cost-diagnostics} records the reporting-budget
implication rather than claiming a measured scaling result.

These are limits of the simulation evidence. Physical switch/NIC access,
real-GPU execution and deployment are outside the agreed study scope and
are not prerequisites for completing the joint simulation.

\section{Conclusion}
Fast confirmation of communication impairment is the primary constraint in
LIMER's active-active recovery study. Class-separated observations,
representative probes and causal traffic compensation meet the registered
per-condition empirical below-100-ms target in a declared 10-Gb/s one-hop
model: 400/400 timely faults, largest per-condition p95 75.4~ms and no
alarms in the tested controls. Same-report contrasts and retained failures
show that observation opportunity, fault treatment and probe cost determine
the scope of that result. Separate simulations demonstrate when a local
switch helps or hurts and check versioned activation with launched-work
ownership. Integrating the primary detector and its costs into that route
loop, then completing the matched scheduling/route comparison, remains
necessary before claiming a robust simulated recovery system.

\bibliographystyle{ACM-Reference-Format}
\bibliography{refs}
\appendix

\section{Alternative Observation Designs}
\label{app:observers}
This appendix records earlier observation components used as references or
fallbacks. They are not interchangeable versions of one evaluated joint
system. Their original traffic models and failed gates remain unchanged.

\subsection{Counter, Transport and Hybrid References}
Hard faults use switch-facing link-state notifications. For gray faults,
the prototype samples switch-facing receive counters at a nominal 20-ms
interval, calibrates healthy burst rates and flags degraded bursts or
deep stalls. A worker step gate then compares affected communication to
fault-free per-slot timing. The observer runs in the management process,
not a switch ASIC, and injector qdisc state is excluded from its inputs.

An opt-in Linux per-socket retransmission sampler records cumulative
counters, deltas, connection identity and missing reads on the same time
axis. It is observation-only. Positive retransmission deltas can also
arise from congestion or reordering, and a delayed first delta cannot
establish a 100-ms bound. This passive observation does not validate
the separate hybrid detector or connect it to automatic recovery.


We additionally implement an isolated passive host decision using
extended Linux \texttt{TCP\_INFO}. ACK-based recovery and reordering
feedback motivate this observation choice~\cite{rack2021}; the rule
interprets Linux state rather than implementing RACK-TLP. Every outgoing
A/B connection is sampled continuously at a nominal 5~ms, without
injector timing or location as input. The frozen rule requires at least
two retransmitted segments in a trailing 100-ms window, outstanding
data and 30~ms since initial loss evidence. A new DSACK or reordering
feedback event vetoes a decision for 100~ms. Missing reads, gaps above
15~ms and connection/counter discontinuities clear evidence; repaired
holes alone do not erase recent repeated-loss evidence. Each connection
generation can issue one decision. This identifies a sender/fabric
transport path, rather than a physical hop, and remains disconnected
from automatic rerouting.


An isolated packet-simulation candidate adds two-ended batch counters,
motivated by alternate-marking measurement and its guard requirements~\cite{alternateMarking2022}.
This is a custom observation model, not a reproduction of that standard.
Each 1500-byte wire packet carries a 12-byte batch/sequence header.
Source acceptance and destination delivery snapshots are independently
reported every 10~ms, with a 50-ms receiver guard and explicitly charged
report transport. Only path, batch, side, count, sample time and report
arrival time enter the detector. It pairs complete batches and requires
two missing deliveries within a trailing 100-ms report-acceptance window.
Two active zero-deficit batches clear history, including before
confirmation, and rearm the rule; idle batches do not prove recovery.
Missing or invalid batches interrupt evidence, and reports older than
200~ms after batch closure are unusable. Decisions include report delay
and are produced by chronological input replay. This observes a delivery
deficit under the guard, without proving permanent loss or its cause.
It remains disconnected from a dynamic route controller.

The hybrid candidate retains this mature lane and adds continuous
64-byte wire probes on the same UDP socket, queue, path and receiver as
data. Probes run at 35,000 packets/s/path independently of fault time,
including during data pauses. Size-independent IID packet loss affects
both traffic types. A separate early lane samples cumulative accepted
counts and delivered counts simultaneously at closure plus 10~ms. It compares
adjacent complete pairs and alarms after two consecutive windows with at
least 32 accepted packets and delivered-count increments no greater than
60\% of accepted-count increments. Gaps, invalid or unavailable pairs and
nonsevere windows break the evidence. Both cumulative samples are monotonic and share the same timestamp;
RX credit cannot exceed simultaneous accepted traffic.
This early signal concerns delivery service, not identified permanent
loss. The observable guard identifies the lane; fault truth remains
scorer-only. The combined rule retains the first alarm per path in each
lifecycle; multiple-episode rearming is not validated. Both lanes execute
natively on delivered-report events, with independent Python replay.

\section{Additional Detection Evidence and Development Records}
\label{app:detection-evidence}
The following records retain prior experiments and negative results.
The primary candidate and validation are in Sections~\ref{sec:detection}
and~\ref{sec:primary-validation}. Cross-model counts below are not pooled
into a detector ranking or into the primary 400-fault denominator.

\subsection{Transport Diagnostics and Earlier Validation}
An observation-only TCP retransmission pilot runs three independent
lifecycles at each configured loss level (0.5, 1, 2, 5, 10 and 25\%),
three short healthy references, and three repetitions each of matched
nominal, paced-B-congestion, round-pause and A-reordering controls.
All 33 attempts remain in the ledger: 30 complete, whereas the three
25\% loss attempts exceed the 120-s round wait. Across 615,998 socket
samples, every counter read is available. Availability does not imply
that the signal identifies the cause of retransmission.

For the affected rank-2 A connection, the first positive counter delta
arrives 44.5--1158.1~ms after 0.5\% loss injection, 59.4--157.4~ms at
1\%, 59.9--152.6~ms at 2\%, and 47.7--62.7~ms at 5\%. These are ranges
of three raw signal observations per level, not detector decision
latencies or confidence bounds. They already preclude a uniform
100-ms claim for the present observer.

Neither nominal nor round-pause controls produce a positive delta;
each congestion control produces at least one delta somewhere in the
observed sockets, and each reordering control produces 730--892
retransmissions on rank-2 A. Congestion traffic is paced at 50~Mb/s,
with measured sender rates of 49.95--49.97~Mb/s. These controls configure
no random packet loss; queue drops are not independently ruled out.
The existing rate-counter detector emits at least one suspect in all
12 matched controls, which is not a final host-confirmed recovery alarm.
The pilot observer remained disconnected from the controller. Its raw
signals motivate the independently validated host candidate below; they
are not reclassified retrospectively as successful decisions.

\paragraph{Packet-granularity diagnostics.}
Eight exploratory capture runs compare original veth settings with
TSO/GSO/GRO disabled, once each at 0, 0.5, 1 and 2\% configured loss.
All complete with zero kernel capture drops and windows covering the
first 100~ms after injection. Original settings expose IP units as large
as 37,700 bytes; disabling offloads bounds observed IP lengths at
1500 bytes but still yields a 1253.7-ms first receive sequence gap in
the 0.5\% case. Disabling device offloads alone consequently does not
resolve the latency constraint or validate the injector's loss unit.

A distinct four-run diagnostic moves loss injection to an ingress IFB
after sender segmentation, retaining the shaping and disabling offloads.
The healthy run shows no sequence gap or positive retransmission delta;
one run each at 0.5, 1 and 2\% loss shows first receiver gaps at
10.6, 10.9 and 9.2~ms and first positive socket counter samples at
22.3, 22.4 and 19.7~ms. These are exploratory signal times under a
changed fault placement, with no new controller trigger or held-out
false-alarm validation. Reordering can also create sequence gaps.
The captures cover an eight-second prefix, not the full workload.
Packet-field checks agree with an independent native decoder in four
selected captures; startup failures, incomplete termination statistics
and lossy capture attempts remain separately archived. These diagnostics
motivate prospective testing without replacing the original measurements
or establishing sub-100-ms detection.

\paragraph{Frozen passive host-rule validation.}
Development distinguishes emulator corrections from detector changes.
The first scene produces six nominal false alarms; a separate recipe
adds symmetric 100-Mb/s, 1-ms post-segmentation ingress shaping with
limit 10,000 to all eight sender paths, retains switch-to-worker
shaping, and disables veth offloads. The changed model is not an
algorithm-only gain. A continuous-loss rule also misses promptly
repaired low-rate losses, motivating the frozen recent-evidence rule.
These exploratory trials and failed setup/reporting attempts remain
archived separately from the prospective matrix.

The independent matrix registers 20 lifecycles each at 0.5, 1, 2 and
5\% loss plus five each nominal, paced-B-congestion, round-pause and
10-ms A-reordering controls. Four ranks exchange framed 2-MiB chunks
over one warmup, two baseline and four post-event rounds, with recovery
off. Source, wrapper, container image and configurations are frozen.
Kernel loss draws are uncontrolled and unpaired; seed 20261010
shuffles launch order only. All outgoing connections are monitored
before and after injection. Offline scoring counts the intended rank-2
A path; the detector receives neither its identity nor injection time.

\begin{table}[t]
\caption{Independent TCP host decisions under the symmetric ingress recipe.
All registered fault attempts remain in the timely denominator. P95 is the
nearest-rank delay with nondetections assigned infinity.}
\label{tab:tcp-host-validation}
\centering
\small
\begin{tabular}{@{}lrrr@{}}
\toprule
Loss (\%) & $<100$ ms / trials & Misses & P95 (ms) \\
\midrule
0.5 & 14/20 & 0 & 408.2 \\
1 & 19/20 & 0 & 85.2 \\
2 & 19/20 & 0 & 60.9 \\
5 & 20/20 & 0 & 49.0 \\
\bottomrule
\end{tabular}
\end{table}

100 of 100 registered lifecycles complete with framed transport
integrity. The actual host decision, measured from injection completion,
is below 100~ms in 72 of 80 fault attempts; 8 are late and
0 have no correctly localized post-injection decision. P95 and
timely fractions are both retained: a p95 pass does not mean every
trial passes. The original all-trial deadline gate is FAIL.
There are 0 other or pre-injection alarms in fault trials.
The 20 primary controls produce 0 false alarm events over
178.4~s of observed lifecycle time; all pre-event and unaffected-path
alarms count, rather than only the known faulted connection.

Five separately registered 40-ms reordering stress lifecycles produce
3 false alarm events in 3 of five runs.
These no-random-loss failures remain visible and block a broad causal
fault-detection or automatic-recovery claim. A sender/fabric decision
cannot localize the physical hop, cover an idle path or establish a
GPU/RDMA hardware guarantee. The primary gray-fault constraint remains
unresolved despite the bounded timing measurements.


\paragraph{Two-ended packet-simulation validation.}
Eight independent one-hop UDP streams use four rank labels and two
plane labels; they do not execute a collective or a training topology.
Low and high profiles offer 80~Mb/s on 100-Mb/s links and 1~Gb/s on
10-Gb/s links, respectively, with 1-ms propagation. Packets are paced
from 1 to 250~ms; observations continue through 600~ms. Offered rates
include all 1500 wire bytes, including the marking and the pinned
backend's 14-byte PPP header. Receiver snapshots wait 50~ms after
batch closure. A separate modeled FIFO charges 64 bytes per report,
1~Mb/s service and 5-ms propagation; it is not a measured control
network and shares no congestion with the monitored data paths.
Shared clocks and complete local counters are explicit assumptions.

After eight development trials, we freeze the rule, parameters,
binary and simulator libraries before 200 fresh lifecycles: 20 each
at four loss rates in both profiles, plus five each nominal,
below-capacity congestion, source-pause and 40-ms reorder controls per
profile. Fault or pseudo-event phase is independently registered in
50--60~ms. The seed is 20261010, with separate assigned streams and
unique validation runs. All eight paths are observed from startup;
fault timing, identity and private drop counters enter only scoring.
Nominal repeats are deterministic checks, not independent evidence
for a stochastic false-alarm distribution.

\begin{table}[t]
\caption{Frozen two-ended observation rule in isolated packet simulation.
All registered loss attempts remain in the denominator; nondetections
remain in nearest-rank p95 as infinity. Rates are offered wire Mb/s.}
\label{tab:batch-validation}
\centering
\small
\begin{tabular}{@{}rrrrrr@{}}
\toprule
Rate & Loss (\%) & $<100$ ms & Late & Miss & P95 (ms) \\
\midrule
80 & 0.5 & 7/20 & 12 & 1 & 228.2 \\
80 & 1 & 12/20 & 8 & 0 & 138.2 \\
80 & 2 & 20/20 & 0 & 0 & 88.0 \\
80 & 5 & 20/20 & 0 & 0 & 78.6 \\
1000 & 0.5 & 20/20 & 0 & 0 & 72.5 \\
1000 & 1 & 20/20 & 0 & 0 & 68.0 \\
1000 & 2 & 20/20 & 0 & 0 & 69.2 \\
1000 & 5 & 20/20 & 0 & 0 & 68.1 \\
\bottomrule
\end{tabular}
\end{table}

All 200 primary lifecycles pass final packet accounting. The rule
produces 139 of 160 correctly localized loss decisions below 100~ms;
the all-trial deadline gate is FAIL. The 40 primary controls produce
0 false alarm events over 24 simulated lifecycle seconds
(192 path-seconds). Fault trials produce 0 pre-event or
unaffected-path alarms. These simulated times are not pooled with
Linux wall-clock measurements, and the different UDP workload is
not a matched comparison establishing improvement over the TCP rule.

All 25 separately registered stress lifecycles drain with audited
packet accounting. Under the original loss-only scoring, 80-ms
reordering and 160\%-capacity competing load each trigger alarms in
five of five no-permanent-loss controls. Sparse 1-Mb/s data yields
zero timely decisions and five nondetections at 0.5\% loss. Idle
traffic yields zero timely, one late and four missed decisions;
10\% report loss yields zero timely, two late and three missed
decisions, with 361 of 2000 scheduled batch pairs unavailable.
An unavailable report, including both reports missing, is never
interpreted as zero received data. A packet arriving beyond the
guard is indistinguishable from a permanently lost packet at the
decision deadline; the original control failures remain visible.

Each lifecycle schedules 800 reports, or 51.2~KB of modeled control
traffic. During steady reporting, the eight paths offer 819.2~Kb/s,
81.92\% of the collector service rate; marking uses 0.8\% of each
wire packet. These are modeled byte and service costs. Exact replay
recomputes all 179620 archived reports, decisions and scores,
checking source, binary, library and input hashes. Final UDP packet
conservation does not establish numerical reduction or coordinated
cutover correctness. An initial 1512-byte wire setup failure and its
seven unrun development registrations are preserved; the corrected
eight-run development ledger precedes the fresh validation.

\paragraph{Frozen hybrid delivery detection.}
The archived failures above motivate a new prospective packet-model
campaign. Development v1 retained sparse-loss and overload deadline
failures; v2's early age-deficit lane alarmed in a mild-congestion control.
V3 replaced that lane with delivery service rate, but its first independent
validation false-alarmed on a high-throughput source pause. V4 aligns TX/RX
sampling times and rejects offset samples before fresh development and
validation. All failed ledgers remain available. A separate stale-binary setup
failure retains eight failed and 129 unrun attempts before corrected compilation. The frozen validation has
740 fresh lifecycles: 60 each at four loss rates in both profiles and in
three weak-loss stress conditions, 20 each for severe delay and overload,
and 40 negative controls. Eligible slowdowns are prospectively defined
as 80-ms additional delivery delay on 25\% of packets or 80\%-capacity
additional competing traffic. Loss, delay and overload remain separate
conditions; the slowdown labels do not rewrite M8's loss-only failures.
Run IDs 41000+ and hidden event phase in 50--60~ms are registered before
execution. All eight paths are monitored; the injected path is fixed at
2/A. The detector sees neither injection time nor path truth.

\begin{table}[t]
\caption{Hybrid packet-model validation. Labels $r/p$ give data wire
Mb/s and loss percent. Every fault attempt remains in its condition's
denominator; misses remain in p95 as infinity. Sparse, idle and report-loss
conditions have 0.5\% packet loss.}
\label{tab:hybrid-validation}
\centering
\small
\setlength{\tabcolsep}{3pt}
\begin{tabular}{@{}lrrrr@{}}
\toprule
Condition & $<100$ ms & Late & Miss & P95 ms \\
\midrule
80/0.5 & 60/60 & 0 & 0 & 78.5 \\
80/1 & 60/60 & 0 & 0 & 70.5 \\
80/2 & 60/60 & 0 & 0 & 64.8 \\
80/5 & 60/60 & 0 & 0 & 62.8 \\
1000/0.5 & 60/60 & 0 & 0 & 66.8 \\
1000/1 & 60/60 & 0 & 0 & 64.0 \\
1000/2 & 60/60 & 0 & 0 & 63.3 \\
1000/5 & 60/60 & 0 & 0 & 62.4 \\
Sparse 1 Mb/s & 60/60 & 0 & 0 & 79.5 \\
Data idle & 59/60 & 1 & 0 & 80.7 \\
Reports 10\% loss & 59/60 & 1 & 0 & 77.7 \\
Delay 80 ms & 20/20 & 0 & 0 & 61.8 \\
Overload & 20/20 & 0 & 0 & 34.0 \\
\bottomrule
\end{tabular}
\end{table}

All 740 lifecycles pass packet and report accounting with exact
native-versus-Python alarm agreement. The hybrid rule yields
698/700 correctly localized fault decisions
below 100~ms; every registered eligible condition meets the 95\%
empirical target. The largest per-condition p95 is 80.7~ms.
The 40 controls produce zero alarms, and fault trials produce zero
pre-event or unaffected-path alarms. Deterministic nominal repetitions
are implementation checks, not a stochastic false-alarm distribution.
All scheduled observations, including wholly missing report pairs, are
accounted for. The separate all-trial/no-unavailable-report gate remains
FAIL; it is not replaced by the empirical fraction rule.
Wilson two-sided 95\% intervals accompany the released ledger:
60/60 has a lower bound of 93.98\%, and 20/20 has 83.89\%.
Thus an empirical pass does not prove a population lower bound of 95\%.

The result incurs substantial observation cost: 17.92~Mb/s of
probes per path, or 17.92\% of a 100-Mb/s link. With 80-Mb/s data,
the nominal total is 97.92\% even without competing traffic. The
10\%-capacity background control therefore offers mild overload rather
than below-capacity traffic; it still produces no severe-impairment alarm.
Early and mature samples each receive three independently lost 64-byte
copies on a separately serialized 10-Mb/s, 1-ms collector. All 4800
copies are charged, including lost copies: 307.2~KB/lifecycle and
4.9152~Mb/s steady load, 49.152\% of collector capacity. The workload
stops at 250~ms and observations drain through 600~ms. This collector
does not share data-plane congestion. The gain over M8 combines changed
observation traffic, report service and rules; it is not a matched
algorithm-only ablation.

Fifteen fresh boundary trials preserve failures outside the primary
envelope. No-probe idle traffic gives
0/5 timely decisions;
120-ms report propagation gives
0/5, and size-selective loss
affecting only packets over 128 bytes gives
0/5.
The delayed-report boundary observes until 800~ms to drain its tail.
In the sparse selective boundary, 4/5 lifecycles realize no
business-packet damage; these configured-fault outcomes are not that many
demonstrated failures to detect actual damage. A separate active-load
supplement uses 20 matched seeds at 80~Mb/s. IID and selective arms have
identical business send/loss counts, 123 dropped
business packets in each arm, and 0
zero-damage pairs. IID gives 20/20 timely alarms;
selective gives 5/20 while all probes remain
healthy. All 40 lifecycles pass accounting and exact native replay. This
matched contrast demonstrates an actual probe-representativeness limitation,
not an assumption of uniform fault coverage. Same path alone is insufficient.
These are one-hop UDP observations with a fixed
fault location, not an RDMA collective or joint route-control result.

\paragraph{Representative business-delivery observations.}
The size-selective failure motivates class-separated counters and receiver
gaps, without discarding the earlier results. Corrected 100-Mb/s development
has 16/20 timely selective decisions and 13/20 application-loss decisions,
with zero alarms in 65 controls. The same-report shadow gives 7/20 and
4/20. One application-loss registration realizes no business drop; it
remains a miss in the denominator. Lowering the 40-ms gap guard retrospectively
false-alarms in all ten lifecycles of each dense, sparse and application
40-ms delay control. The initial setup also incorrectly injected 80~ms
under a sparse-40-ms label: its 19 completed, five interrupted and 125 unrun
registrations remain archived and supply no acceptance evidence.

Changing probe shape at the old nominal 17.92-Mb/s byte budget improves
selective development to 20/20, but leaves application loss at 18/20 and
idle loss at 2/5. Smaller probes on fresh, unpaired draws give 17/20,
14/20 and 5/5, respectively. Both arms retain full-coverage FAIL and
all 60 controls remain alarm-free. Actual packet-boundary costs are
557,760 versus 558,000 probe bytes/path; this difference is charged.

We therefore register an explicitly different envelope: 10-Gb/s links,
1500-byte business-shaped probes and the unchanged 40-ms detector guard.
Constant 35,000-pps probes cost 420~Mb/s/path (4.2\% of link capacity),
or 525\% and 42\% of business bytes at 80~Mb/s and 1~Gb/s. Following
70/70 timely development faults and 30 alarm-free controls, the unchanged
source passes fresh 460-case validation with 400/400 timely faults and
60 alarm-free controls (Table~\ref{tab:business-validation}). The
same-report shadow baseline gives 58/60 in each 80-Mb/s selective and
application-loss condition, and all other fault decisions timely: it also
passes every condition's 95\% target. The improvement largely comes from
observation opportunity, not an established novel decision rule.

\begin{table}[t]
\caption{Constant representative-probe validation on 10-Gb/s links.
Weak faults use 0.5\% loss; delays affect 25\% of packets. All attempts
remain in the denominator, with zero late or missed decisions.}
\label{tab:business-validation}
\centering
\small
\setlength{\tabcolsep}{4pt}
\begin{tabular}{@{}lrr@{}}
\toprule
Condition & $<100$ ms & P95 ms \\
\midrule
Selective, 80 Mb/s & 60/60 & 63.2 \\
Application loss, 80 Mb/s & 60/60 & 61.2 \\
Selective, 1 Mb/s & 60/60 & 64.9 \\
Business idle & 60/60 & 62.6 \\
Selective, 1 Gb/s & 60/60 & 56.1 \\
Application loss, 1 Gb/s & 60/60 & 55.9 \\
Network delay 80 ms & 20/20 & 53.0 \\
Application delay 80 ms & 20/20 & 54.6 \\
\bottomrule
\end{tabular}
\end{table}

\section{OptCC Reconstruction and Artifact Checks}
\label{app:optcc}
Section~\ref{sec:baseline-reproduction} reports the matched schedule timings.
The following original artifact and implementation records establish how
that independent comparator was constructed.

\paragraph{Artifact validation completed.}
We archived the six-file public reference artifact and compiled its three
standalone C++17 schedule generators~\cite{optccArtifact2026}. Four example
cases completed and produced parseable DAGs encoding flow dependencies
(Table~\ref{tab:optcc-artifact}). The single- and multiple-straggler
examples passed the generators' internal validation. The multi-GPU example
matched the driver's reference makespan; this is a model-level check,
not a measured network completion time. The latter example also assumes
NVLink-rich servers and is retained only as an artifact check.

\begin{table}[t]
  \caption{Completed OptCC artifact checks. Flow counts are generated DAG
  entries, not GPU execution results or reproduced performance measurements.}
  \label{tab:optcc-artifact}
  \centering
  \small
  \begin{tabular}{@{}lrl@{}}
    \toprule
    Reference case & Flows & Check \\
    \midrule
    Single straggler, $L=2$ & 132 & PASS \\
    Single straggler, $L=1.5$ & 157 & PASS \\
    Two stragglers, $L=(2,1.5)$ & 274 & PASS \\
    Four GPUs/server, $L=2$ & 198 & Ref. match \\
    \bottomrule
  \end{tabular}
\end{table}

\paragraph{Independent packet-level reconstruction.}
The simulator integration and experiment driver are absent from the
public snapshot, preventing an exact reproduction of a published
performance figure. We pinned SimAI commit \texttt{f5efb5a} and
its ns-3 backend \texttt{7e3cb5b}. The original simulator's eight-GPU
NVLink microbenchmark completed before the graph importer was added.
For the single-GPU-per-server setting, our standalone replay uses SimAI's
network frontend and RDMA packet backend without the native collective
executor's memory or reduction costs. Both OptCC and the ordinary
reduce-scatter/all-gather ring use this same network-only runner.

We reconstruct payload sizes from the disclosed section-size formulas;
the JSON itself contains dependencies but no tensor or message bytes.
Both schedules transmit $2(p-1)N$ data bytes in aggregate in the tested
$L=2$ cases. OptCC synchronization nodes carry an additional 20-byte
control payload each. This validates byte accounting and schedule
execution, not numerical AllReduce correctness. All receives are posted
before roots launch, and each dependency is released at its source-side
send completion or destination-side ACK-complete receive notification.

Transport gates exposed missing or premature whole-flow send callbacks
at packet-size boundaries. We repaired final-packet serialization
notification in the shared backend and repeated all 14 transport gates.
The replay uses chronological serial event processing: the original
multithreaded scheduler drained the OptCC graph but violated dependency
timestamps even with one worker thread. Three scheduler micrographs,
the 132-node $L=2$ graph (255 edges), the 157-node $L=1$ bubble graph,
and a ring smoke test subsequently passed completion and trace checks.
The earlier failed attempts are retained.

\paragraph{Additional deterministic checks.}
Within this declared setting, the OptCC schedule reduces network
makespan by 10.4--28.8\% relative to the matched ring. Four healthy-ring
references and two preregistered segmentation checks also completed,
bringing the measurement denominator to 30/30. At $p=16$, $N=15$~MiB,
increasing $k$ from 4 to 8 or 16 changed makespan from 1.986~ms to
2.143~ms or 2.886~ms; the primary $k=4$ arm is retained. These measurements
establish an executable comparator under bandwidth asymmetry. They do
not measure route switching, packet-loss detection, recovery overhead,
or end-to-end training throughput.

\section{Execution and CPU Emulation Records}
\label{app:execution}
The framed-TCP and native pending-flow targets have different execution
contracts. These records retain their original integrity checks, failures
and timing scopes; none proves exactly-once numerical AllReduce.

\subsection{Framed-TCP Ownership Rules}
Frames identify the round, step and route version and carry a payload
length and checksum. Workers reject invalid frames and discard already
consumed duplicates. The coordinator changes only the affected route
slots, while interrupted steps retain completed send or receive work
where the route is unchanged.

An in-flight receive has one owner per round and step. If a send-side
abort interrupts the step while the receive route is unchanged, redo
reuses the existing reader, including its partial frame and result.
Creating another reader could consume the expected frame in an abandoned
attempt. If the receive route changes, the superseded reader must stop
before replacement. Deterministic tests cover consumed frames and
partial headers across this abort; topology runs separately check
completion, plan agreement and integrity. These checks do not constitute
an exactly-once numerical AllReduce proof.

\subsection{Recovery Outcomes and Native Transaction Cases}
Fresh hard-fault, gray rate-degradation and transient checks each
complete three lifecycles, with framed-transport integrity, schedule
conformance and route-version agreement. All three hard-fault and all
three sustained gray-degradation runs commit a localized route plan.
Only two of the three transient runs suppress recovery; the remaining
transient run commits after 1086.5~ms. Completion and transport integrity
therefore do not establish correct transient discrimination. These nine
runs also do not establish reliability for random loss.

The first original live-policy switch and stay attempts at 25\% loss
both select stay with unknown alternate spare; both subsequently exceed
the 120-s round wait. The other 27 queued repetitions are explicitly
unrun after the branch-gate failure. The guarded switch pilot completes
11 collective rounds (44 worker rounds), with only rank 2's three A
slots changed and zero checksum or version errors. The guarded stay
pilot selects its expected branch but also times out. Its failure blocks
expansion to a repeated reliability claim. A policy label alone is not
evidence that the selected arm minimizes completion cost.

\paragraph{Scripted pending-flow simulation gate.}
A separate native SimAI target tests runtime route changes before online
detector integration. At launch it selects the committed plane, posts one
matching receive and retains that transfer's ports and version. A shared,
reliable modeled management FIFO charges 64 bytes per PREPARE, READY and
COMMIT, with 1-Gb/s serialization and 1-$\mu$s propagation by default.
Activation follows readiness and commit delivery at all four ranks.
Management messages are not packetized on the RDMA data network.

All nine prospectively registered engineering cases pass: eight workloads
complete with transport-map drain, and one deliberately nonconsecutive
version is rejected. Independent auditing reconstructs all 582 trace
records and verifies frozen inputs, source, binary and 39 shared libraries.
The default transaction requested at 5~$\mu$s activates at
13.145~$\mu$s; two pending affected ring transfers move while launched
transfers retain ownership. Two consecutive versions move four pending
transfers. With 250-$\mu$s propagation, activation occurs at
759.609~$\mu$s after workload completion at 221.994~$\mu$s, so no
transfer moves. A missing READY prevents activation and reaches its
registered timeout at 1005~$\mu$s. Source-side and destination-side
dependency chains and an unequal-size join also drain correctly. The
equal-rate, impairment-free ring completes in 221.994~$\mu$s without
control and 222.678~$\mu$s with a switch. These scripted checks establish
pending-flow execution under the modeled protocol, not detection efficacy,
beneficial recovery, cancellation/replay or numerical reduction. A
severely stalled launched transfer cannot be rescued by this mode.

\section{Reference Policies and ACK-Progress Development}
\label{app:online}
The original headroom heuristic below belongs to the CPU prototype.
The remaining-work heuristic in Section~\ref{sec:recovery-design} belongs
to the separate native online loop. Their results are not interchangeable.

\subsection{Earlier Headroom Policy}
The current heuristic compares the affected plane's calibrated demand
with the alternate plane's estimated spare service rate. A burst median
and busy fraction provide demand and spare proxies; the rule switches
when spare multiplied by a fixed reuse factor of 1.25 reaches the
pre-fault demand. The factor is retained rather than tuned from the new
outcomes. This rule is a heuristic, not a fitted model of total episode
completion cost.

The original live estimate becomes unknown when its 256-sample window
contains no bursts. A fault-induced collective pause can cause that
condition even on a reachable alternate link. An opt-in guarded variant
keeps the healthy spare estimate as a ceiling, takes a lower known live
estimate, and uses the healthy reference when the live estimate is
unknown. An unknown healthy reference still selects stay. This variant
assumes stationary alternate capacity and cannot detect a silent capacity
decrease on an idle port. Both inputs and the assumption are logged.

\paragraph{Native causal ACK-progress development.}
A separate native target samples sender acknowledgment progress every
2~ms, combining completed flow bytes with active queue-pair ACK credit.
The rule calibrates the median of four positive fully busy service
samples, then alarms after three busy samples below half that reference.
Idle windows interrupt evidence. Reports occupy 64 bytes each on a
separate reliable 10-Mb/s FIFO with 1-ms propagation. The detector runs
on delivered reports inside simulator callbacks; injection metadata is
excluded. This signal concerns ACK progress, not permanent-loss proof.

Four deterministic development scenes execute three consecutive four-rank
3-MiB rings with 1-Gb/s accesses. Persistent rank-2 A service reduction
to 0.2~Gb/s at 50~ms alarms after 9.1536~ms; healthy and fully restored
2-ms reduction scenes produce no alarms. Changing report propagation to
120~ms delays the alarm to 128.1536~ms, failing the deadline. An earlier
device-only restoration left a newly created queue pair capped at the
reduced rate for its 31.563-ms transfer and produced a registered control
alarm. That ledger remains intact. The amended injector changes directed
configured NIC service and active DCQCN pacing caps at both reduction and
restoration; detector logic is unchanged. This model correction is not an
algorithm-only gain or a hidden physical-link fault simulation. Independent
replay verifies all 2920 delivered reports and 226492416 ACK bytes in
the amended four-case ledger. These are development cases, not broad
held-out validation of the primary constraint.

\section{Additional Cost and Scaling Diagnostics}
\label{app:cost-diagnostics}
\subsection{Static Calibration Setup}
The healthy 200/200 reference completes in 0.222~ms, and a 200/20
reference completes in 1.128~ms. At 40/200, switching reduces makespan
by 22.9\%; at 40/20 it increases makespan by 33.0\%. These are static
network comparisons, not LIMER controller results. Three initial
40/20-Gb/s attempts aborted before simulated time advanced because the
default ECN tables omitted those link speeds. We retained them, added
explicit thresholds scaled from the supplied 25-Gb/s row and repeated
all six calibrations; the two constrained-B references were separately
registered. Runtime delay, background demand and finite-horizon costs
remain necessary for the route-choice study.

\subsection{Severe-Loss Fixed-Route Contrast}
A separate fixed-route diagnostic retains the framed-TCP workload:
25\% random loss on rank-2 A, 10-Mb/s rank-2 B access, and eight
post-fault rounds. All three fixed-stay attempts exceed the 120-s
round wait; they complete zero, one and zero post-fault rounds,
respectively. All three immediate safe-boundary localized-switch
attempts complete all eight rounds with integrity and locality checks,
in 84.804--85.029~s from fault injection to the last round's completion.
The switch changes only rank 2's three A slots. Source and container
image are shared, but realized packet-loss draws are not paired.

This diagnostic omits detection delay and in-flight cutover by moving
at the next safe round. It cannot establish adaptive LIMER superiority
or quantify the counterfactual completion time of a timed-out arm.
It does show why insufficient estimated spare is not an optimality
label: staying on a highly impaired TCP path can fail to finish even
when the alternative is slow. Combined with the static calibration's
opposite preferred arms, the result motivates a finite-episode cost
objective that distinguishes bandwidth reduction from random loss.

\subsection{Scaling with SimAI}
The current baseline includes 5- and 16-rank single-plane graphs; the
dual-plane calibration includes four logical ranks. The planned crossed
schedule/route matrix and expansion up to 128 ranks remain unrun.
At unchanged 10-ms two-ended reporting, the modeled collector load is
204.8~Kb/s per logical rank: eight ranks exceed the current 1-Mb/s
service rate, and 128 ranks would offer 26.2144~Mb/s. This is an
algebraic budget implication, not a measured 128-rank result. Simulated
control delays and actual implementation latencies remain separate.

\section{Implementation Environments}
\label{app:implementation}
The following environment and test counts describe their recorded execution
checkpoint, not a newly run verification or a current full-suite count.

The CPU harness uses Mininet, kernel Open vSwitch and a four-worker
framed TCP ring. Each worker has a persistent outgoing and incoming
connection on each plane. Faults act on a directed worker egress, while
the detector observes the switch-facing peer. Workers stream step events
and plan versions to the coordinator; results retain per-round byte
counts, fingerprints, checksums, route changes and failures.

Fresh execution uses isolated rootless Docker containers on an approved
Linux host, Ubuntu 24.04/Python 3.12 and a nominal 4-CPU/4-GiB budget per
lifecycle, with at most two concurrent Mininet runs. Namespace and sysfs
mounting require privileged mode inside the rootless daemon. Resource
limit warnings and actual sampling intervals are retained. The complete
Linux unit suite passes 292 tests on the Linux Python 3.12 host. These checks establish implementation
behavior, not deployment on GPU collectives or switch hardware.

The separate SimAI adapter uses the pinned network frontend and packet
backend. The single-plane comparator reconstructs graph byte sizes and
endpoint dependency releases. The dual-plane calibration adds two
physical network endpoints per logical rank, with disjoint switches and
plane-consistent sends and receives. It changes static slot placement before launch and includes no runtime
detector or cutover delay. A separate dynamic target registers a single
matching receive at each launch and retains its physical endpoints and
version through completion. Only pending affected transfers can change
planes; launched transfers remain on their original route.


\end{document}
